\section{سیگنال‌ها در لینوکس}

\subsection*{بخش آ – رفتارهای پیش‌فرض پس از ورود سیگنال}

برای هر سیگنال فقط پنج رفتار پیش‌فرض ممکن است (راهنمای \code{signal(7)}):
\begin{itemize}
  \item \code{Term}: پایان‌دادن به پردازه (مثل \code{SIGTERM}، \code{SIGINT}، \code{SIGKILL}، \code{SIGALRM}).
  \item \code{Core}: پایان‌دادن به پردازه و ساختن فایل \lr{core dump} از حافظهٔ پردازه (مثل \code{SIGSEGV}، \code{SIGABRT}، \code{SIGFPE}).
  \item \code{Ign}: نادیده‌گرفتن سیگنال (مثل \code{SIGCHLD}، \code{SIGURG}).
  \item \code{Stop}: متوقف‌کردن پردازه (مثل \code{SIGSTOP}، \code{SIGTSTP}).
  \item \code{Cont}: ادامه‌دادن پردازهٔ متوقف‌شده (\code{SIGCONT}).
\end{itemize}

دو مورد را توضیح می‌دهیم:
\begin{itemize}
  \item \textbf{\code{Term} (مثلاً \code{SIGTERM}، شمارهٔ ۱۵):} هسته پردازه را بدون ذخیرهٔ وضعیت متوقف می‌کند و منابعش را (حافظه، فایل‌های باز) آزاد می‌کند. والد با \code{wait} وضعیت پایان را می‌بیند. اگر برنامه گردانندهٔ سیگنال نصب کرده باشد، به‌جای رفتار پیش‌فرض آن اجرا می‌شود (برای مثال برای بستن مرتب فایل‌ها). \code{SIGKILL} و \code{SIGSTOP} را نمی‌توان گرفت، مسدود یا نادیده گرفت.
  \item \textbf{\code{Core} (مثلاً \code{SIGSEGV}، شمارهٔ ۱۱):} وقتی پردازه به آدرس نامعتبر دسترسی می‌دهد، سخت‌افزار تله ایجاد می‌کند، هسته آن را به سیگنال تبدیل می‌کند و پردازه پایان می‌یابد و تصویر حافظه و ثبات‌ها در فایل \lr{core dump} برای اشکال‌زدایی ذخیره می‌شود (اسلاید فصل ۲: «\lr{Failure of an application can generate core dump file}»).
\end{itemize}
\src{تکمیلی (\code{signal(7)})، اشاره به \lr{core dump} و تله در اسلایدها.}

\subsection*{بخش ب – مؤلفه‌های مرتبط با سیگنال در \code{/proc/PID/status}}

این فایل چند فیلد دربارهٔ سیگنال‌ها دارد که همگی \textbf{نقاب بیتی شانزده‌دهی} هستند: بیت شمارهٔ \lr{n-1} روشن یعنی سیگنال شمارهٔ \lr{n} در آن مجموعه است. سه مؤلفهٔ اصلی:
\begin{itemize}
  \item \code{SigBlk} (\lr{Blocked}): سیگنال‌های \textbf{مسدودشده} (نقاب سیگنال). تحویل سیگنال‌های این مجموعه به تعویق می‌افتد تا نقاب برداشته شود. مثال: \code{SigBlk: 0000000000010000} یعنی بیت ۱۶ (سیگنال ۱۷، یعنی \code{SIGCHLD}) مسدود است.
  \item \code{SigIgn} (\lr{Ignored}): سیگنال‌هایی که پردازه \textbf{نادیده می‌گیرد} (رفتار \code{SIG\_IGN}).
  \item \code{SigCgt} (\lr{Caught}): سیگنال‌هایی که پردازه برایشان \textbf{گردانندهٔ سفارشی} ثبت کرده است (مثل \code{sigHandler()} در شکل ۱ صورت تمرین).
\end{itemize}

فیلدهای مرتبط دیگر: \code{SigPnd} و \code{ShdPnd} (سیگنال‌های \textbf{در انتظار} برای نخ و برای کل پردازه) و \code{SigQ} (تعداد سیگنال‌های در صف نسبت به سقف مجاز، مثلاً \code{0/31538}).
\src{تکمیلی (\code{proc(5)}).}

\subsection*{بخش ج – سیگنال‌های تله‌های فراخوانی نامعتبر و خطای ممیز شناور}

\begin{itemize}
  \item \textbf{فراخوانی سیستمی نامعتبر:} \code{SIGSYS} (شمارهٔ ۳۱ در \lr{x86-64}، «\lr{Bad system call}»)، رفتار پیش‌فرض \code{Core}. نکته: در لینوکس یک شمارهٔ فراخوانی ناموجود معمولاً فقط خطای \code{ENOSYS} برمی‌گرداند و سیگنال نمی‌فرستد. \code{SIGSYS} وقتی ارسال می‌شود که یک فیلتر \lr{seccomp} فراخوانی را رد کند (\code{SECCOMP\_RET\_TRAP}) یا در سیستم‌های یونیکس دیگر مثل \lr{BSD} و \lr{Solaris}.
  \item \textbf{خطا در محاسبات (از جمله شناور):} \code{SIGFPE} (شمارهٔ ۸، «\lr{Floating-point exception}»)، رفتار پیش‌فرض \code{Core}. با وجود نامش، علاوه بر خطاهای ممیز شناور (تقسیم بر صفر، سرریز و ...) برای تقسیم صحیح بر صفر یا سرریز صحیح نیز ارسال می‌شود. سازوکار آن همان تله است که اسلایدها توضیح می‌دهند (تقسیم بر صفر یک وقفهٔ نرم‌افزاری ناشی از خطاست): \lr{CPU} استثنا ایجاد می‌کند، هسته آن را به \code{SIGFPE} تبدیل و به پردازه ارسال می‌کند. نوع خطا از \code{si\_code} (مثل \code{FPE\_INTDIV}، \code{FPE\_FLTDIV}) مشخص می‌شود.
\end{itemize}
\src{تکمیلی (\code{signal(7)}، \code{sigaction(2)})؛ مفهوم تله در اسلاید.}

\subsection*{بخش د – چرا امکان ماسک‌کردن وقفه لازم است؟}

\begin{itemize}
  \item \textbf{حفاظت از بخش‌های بحرانی و جلوگیری از شرایط رقابتی:} اگر گرداننده در میانهٔ به‌روزرسانی یک ساختمان داده (مثلاً صف زمان‌بندی) اجرا شود و همان ساختمان داده را تغییر دهد، داده ناسازگار می‌شود. با ماسک‌کردن، اجرای اتمیک تضمین می‌شود (برای نمونه هسته هنگام تغییر ساختمان‌های مشترک با گرداننده وقفه را غیرفعال می‌کند، و برنامه با \code{sigprocmask} سیگنال‌ها را مسدود می‌کند).
  \item \textbf{جلوگیری از تودرتوشدن و بازورود:} هنگام اجرای گرداننده، همان وقفه (یا وقفه‌های با اولویت پایین‌تر) نباید دوباره آن را قطع کند؛ وگرنه سرریز پشته یا ناسازگاری وضعیت رخ می‌دهد. سیستم‌عامل ذخیره و بازیابی وضعیت را هم باید بدون مزاحمت انجام دهد (اسلاید: سیستم‌عامل ثبات‌ها و شمارندهٔ برنامه را ذخیره می‌کند).
  \item \textbf{مرحلهٔ راه‌اندازی و تغییر حالت:} هنگام بوت، پیش از تنظیم جدول وقفه و گردانندها، یا هنگام سوئیچ زمینه، دریافت وقفه خطرناک است.
  \item \textbf{اولویت‌بندی و کنترل بار:} وقفه‌های مهم‌تر روی کم‌اهمیت‌ها مقدم می‌شوند و طوفان وقفه (بار زیاد) را می‌توان مهار کرد. وقفه ازدست‌نرفته است و تا برداشتن نقاب در حالت انتظار می‌ماند (در سیگنال‌های استاندارد چند ارسال پشت‌سرهم یکی حساب می‌شوند). وقفه‌های حیاتی (مثل خطای سخت‌افزاری شدید) غیرقابل‌ماسک (\lr{NMI}) هستند تا هرگز نادیده گرفته نشوند.
\end{itemize}
\src{تکمیلی؛ اشاره به ذخیرهٔ وضعیت هنگام وقفه در اسلاید فصل ۱.}
